\begin{lemma}[Chernoff-Hoeffding bound]
  \label{lemma:chernoff}
  Let $X_1, \ldots, X_m$ be independent random variables in $[0,1]$, and let
  $\mu = ({1}/{m}) \E[ \sum X_i]$.  Then, for all $\eps \ge 0$,
  $$
  \Pr\left[ \frac{1}{m} \sum X_i \ge \mu + \eps \right] \le e^{-2\eps^2 m}\,.
  $$
\end{lemma}

\section{Computing approximate pure Bayes-Nash equilibrium}
\label{sec:bayes}

Bayesian games model the situation where the players' knowledge of the
world is incomplete.  In this paper we focus on Bayesian games
with two players, but the results generalize to an arbitrary number of
players. More details on Bayesian games can be found in most Game
Theory textbooks, for example in~\cite{FudenbergTirole}.

In a \emph{Bayesian game} the payoff of the players depends on the
state of the world in addition to the players' strategies. In a
situation with two players, the row player and the column player, each player is presented with
a signal, called type, about the state of the world $\ta_{\mathrm{row}}\in \Ta_{\mathrm{row}}$
and $\ta_\mathrm{col}\in\Ta_\mathrm{col}$, respectively.  The types are distributed
according to some joint distribution $\cP$ and are not necessarily
independent. The types determine the values in the payoff matrices $M_{\mathrm{row}}(\ta_{\mathrm{row}},\ta_\mathrm{col})$ and $M_\mathrm{col}(\ta_{\mathrm{row}},\ta_\mathrm{col})$, but not their dimensions. Denote 
the set of rows and columns in any of these matrices by $S_{\mathrm{row}}$ and $S_\mathrm{col}$, respectively.
 Each player chooses an action $s_{\mathrm{row}}\in S_{\mathrm{row}}$ and $s_\mathrm{col}\in
S_\mathrm{col}$ from their respective set of actions.
The payoff function of the first player is thus $u_{\mathrm{row}}(s_{\mathrm{row}},s_\mathrm{col},\ta_{\mathrm{row}},\ta_\mathrm{col})=M_{\mathrm{row}}(\ta_{\mathrm{row}},\ta_\mathrm{col})_{s_{\mathrm{row}},s_\mathrm{col}}\in
[0,1]$. The payoff function $u_\mathrm{col}$ is defined similarly.
The payoff matrices, that depend on the players' types, as well as the distribution on types is known to the players
ahead of the game. 

A \emph{pure strategy} for the row player  in a Bayesian game is a function (that by a slight abuse of notation) we denote 
by $s_{\mathrm{row}}:\Ta_{\mathrm{row}}\ra S_{\mathrm{row}}$ that for each type $\ta_{\mathrm{row}}$ as observed by row player associates a strategy $s_{\mathrm{row}}(\ta_{\mathrm{row}})$
that the player chooses to execute. A pure strategy $s_\mathrm{col}:\Ta_\mathrm{col}\ra S_{\mathrm{row}}$ is defined similarly. 

Denote by $\cP_{\ta_{\mathrm{row}}}$ the distribution on the column player's types $\ta_\mathrm{col}$ conditioned on the type $\ta_{\mathrm{row}}$ being observed. 
For a pair of pure strategies $(s_{\mathrm{row}},s_\mathrm{col})$ the payoff function of the row player  is given by 
$$
p_{\mathrm{row}}(\ta_{\mathrm{row}}) = \E_{\ta_\mathrm{col}\sim \cP_{\ta_{\mathrm{row}}}} [ u_{\mathrm{row}} ( s_{\mathrm{row}}(\ta_{\mathrm{row}}), s_\mathrm{col}(\ta_\mathrm{col}), \ta_{\mathrm{row}}, \ta_\mathrm{col})]\,. 
$$

A \emph{pure strategy Nash equilibrium} in a Bayesian game, is a pair of functions $s_{\mathrm{row}}$, $s_\mathrm{col}$ such that for all types
observed, neither player has an incentive to deviate from his current strategy. In other words, for each $\ta_{\mathrm{row}}$, and 
for each $s_{\mathrm{row}}'\in S_{\mathrm{row}}$, 
$$
p_{\mathrm{row}}(\ta_{\mathrm{row}}) \ge \E_{\ta_\mathrm{col}\sim \cP_{\ta_{\mathrm{row}}}} [u_{\mathrm{row}} ( s_{\mathrm{row}}', s_\mathrm{col}(\ta_\mathrm{col}), \ta_{\mathrm{row}}, \ta_\mathrm{col})]\,,
$$
and a similar condition holds for $p_\mathrm{col}$. 

Since a pure Nash equilibrium need not exist in non-Bayesian games, it
need not exist in Bayesian games either. Moreover, while verifying
whether a non-Bayesian two player game has a pure Nash equilibrium is
trivial, verifying whether a pure Bayesian Nash equilibrium exists is
\cclass{NP}-hard~\cite{CS08}. Furthermore, as the example in~\cite{CS08}
demonstrates, this problem remains hard even when the distribution on
types is uniform and the payoff does not depend on the players' types.

A \emph{pure  $\eps$-Bayesian Nash equilibrium ($\eps$-BNE)} is defined similarly to an $\eps$-Nash equilibrium. 
For each observed type $\ta_{\mathrm{row}}$, the incentive to deviate should be bounded by $\eps$:
$$
p_{\mathrm{row}}(\ta_{\mathrm{row}}) \ge \E_{\ta_\mathrm{col}\sim \cP_{\ta_{\mathrm{row}}}} [u_{\mathrm{row}} ( s_{\mathrm{row}}', s_\mathrm{col}(\ta_\mathrm{col}), \ta_{\mathrm{row}}, \ta_\mathrm{col})]-\eps\,.
$$
A similar requirement should hold for the column player. 

We show that for general distributions on types and for some small constant $\eps$, finding a pure $\eps$-BNE 
in games where a pure BNE exists is still \cclass{NP}-hard. On the other hand, we also show that if the distribution on the
players' types is uniform, whenever a pure BNE exists,  a pure $\eps$-BNE can be found in quasi-polynomial time. 

Next we give a quasi-polynomial algorithm for finding an $\eps$-BNE in
Bayesian games with uniformly distributed types.  The algorithm easily
generalizes to the case when the distribution on types is a product
distribution, \ie, when $\ta_{\mathrm{row}}$ is independent from
$\ta_\mathrm{col}$.  To simplify the presentation we further assume
that the type space is of equal size for both players,
\ie, $|\Ta_{\mathrm{row}}|=|\Ta_\mathrm{col}|=k$.  We prove:

\begin{theorem}[Detailed statement of Theorem~1.8]
\label{thm:BNEeasy-formal} In a two-player Bayesian game,
suppose that the types are distributed uniformly on the space $\Ta_{\mathrm{row}}\times\Ta_\mathrm{col}$, and that $|\Ta_{\mathrm{row}}|=|\Ta_\mathrm{col}|=k$, and $|S_{\mathrm{row}}|=|S_\mathrm{col}|=n$. 
Assuming that a pure BNE exists, we can find a pure $\eps$-BNE in time $n^{O((\ln n+\ln k)/\eps^2)}$. 
\end{theorem}

\begin{remark}
  The assumption in Theorem~1.8 that an exact equilibrium exists can be relaxed to a pure $(\eps/2)$-BNE equilibrium existing.
\end{remark}


\begin{proof}
The proof is similar in spirit to the quasi-polynomial $\eps$-Nash algorithm of~\cite{LMM03}, but some additional work 
is needed. We first (approximately) guess the payoffs and the allowed strategies for both players for all possible types. We then use linear programming to produce a (not necessarily pure) $(3\eps/4)$-BNE\@. We then sample from this approximate non-pure BNE to obtain a
pure $\eps$-BNE\@. 

For simplicity, denote $\Ta_{\mathrm{row}}=\Ta_\mathrm{col}=\{1,\ldots,k\}$ and $S_{\mathrm{row}}=S_\mathrm{col} = \{1,\ldots,n\}$. Let $s_{\mathrm{row}}:\Ta_{\mathrm{row}}\ra S_{\mathrm{row}}$, $s_\mathrm{col}:\Ta_\mathrm{col}\ra S_\mathrm{col}$ be 
a pair of pure equilibrium strategies (that we assume exist). Let $p^{\mathrm{row}}_{ij}$ and $p^\mathrm{col}_{ij}$ be the corresponding 
 payoff values. In other words, $p^{\mathrm{row}}_{ij}$ is the expected payoff the row player gets when his type is $i\in\{1,\ldots,k\}$, he plays strategy $j\in\{1,\ldots,n\}$, and the column player plays according to the equilibrium. The equilibrium assumption is that for all $i$,
 $$
 p^{\mathrm{row}}_{i,s_{\mathrm{row}}(i)} = \max_j p^{\mathrm{row}}_{ij},
 $$
 and a similar condition holds for $p^\mathrm{col}$. We first claim that we can recover all the values of $p^{\mathrm{row}}$ and $p^\mathrm{col}$ within 
 an error of $\eps/4$ with probability $>n^{-O((\ln (nk))/\eps^2)}$.
\begin{claim}
\label{cl:BNE1}
There is a randomized polynomial time algorithm that with probability $>n^{-O((\ln (nk))/\eps^2)}$ outputs values $q^{\mathrm{row}}_{ij}$, $q^\mathrm{col}_{ij}$ such that for all $i,j$, 
$|q^{\mathrm{row}}_{ij}-p^{\mathrm{row}}_{ij}|<\eps/8$ and $|q^\mathrm{col}_{ij}-p^\mathrm{col}_{ij}|<\eps/8$.
\end{claim}

\begin{proof}
We show how to approximate $p^{\mathrm{row}}$ well with probability $>n^{-O((\ln (nk))/\eps^2)}$. The claim follows by approximating $p^{\mathrm{row}}$ and $p^\mathrm{col}$ independently.
Pick a uniformly random subset (with repetitions) of $m=\lceil (40\ln (nk))/\eps^2\rceil$ of the column player's types $t_1,\ldots,t_m\in\{1,\ldots,k\}$. For each type $t_r$ guess (uniformly) the column player's strategy $s_r$ on this type. With probability $n^{-m}$ we correctly guess all strategies, \ie, $s_r = s_\mathrm{col}(t_r)$ for all $r$. Set 
$$
q^{\mathrm{row}}_{ij} = \frac{1}{m}\sum_{r=1}^{m} u_{\mathrm{row}}(j,s_r,i,t_r)\,.
$$
In other words, we calculate an estimate of the expected payoff $p^{\mathrm{row}}$ using only the value at the types we've guessed. By the Chernoff-Hoeffding bound \expref{Lemma}{lemma:chernoff}, for each $(i,j)$ the probability 
$$
{\mathbf P}[|q^{\mathrm{row}}_{ij}-p^{\mathrm{row}}_{ij}|\ge \eps/8] \le 2\cdot \exp(-m\eps^2/32) < 1/(4nk)\,.
$$
The claim follows by union bound. 
\end{proof}

Allowing for a blow-up of $n^{O((\ln (nk))/\eps^2)}$ in running time, we may assume from now on that the correct values 
of $q^{\mathrm{row}}$ and $q^\mathrm{col}$ had been computed (up to an error of $\eps/8$), since in the end we can check whether the set of pure strategies obtained 
is a $\eps$-BNE.

Next, we formulate a linear program that obtains a $(3\eps/4)$-BNE with payoff values close to $q^{\mathrm{row}},q^\mathrm{col}$. The variables 
of the program are $X^{\mathrm{row}}_{ij}$ and $X^\mathrm{col}_{ij}$, where $X^{\mathrm{row}}_{ij}$ corresponds to the probability that the row player plays strategy 
$j$ on type $i$ and similarly for $X^{\mathrm{col}}_{ij}$. For each type $i$ let $$M^{\mathrm{row}}_i := \max_j q^{\mathrm{row}}_{ij}$$ be the maximum possible payoff the row player can attain 
on type $i$ under payoffs $q^{\mathrm{row}}_{ij}$. We only allow $X^{\mathrm{row}}_{ij}$ to be non-zero when the corresponding payoff 
$q^{\mathrm{row}}_{ij}> M^{\mathrm{row}}_i - \eps/8$. This guarantees that the solution we obtain is a $(3\eps/4)$-BNE\@. In addition, we enforce
the payoffs to be close to $q^{\mathrm{row}}_{ij}$. For each $i,j$ we have the constraint
$$
\left| q^{\mathrm{row}}_{ij} - \frac{1}{k}\cdot \sum_{y,z} X^\mathrm{col}_{yz} \cdot u_{\mathrm{row}}(j,z,i,y) \right| < \frac{\eps}{4}\,,
$$ 
and a similar condition on the $X^{\mathrm{row}}$'s. 
This linear program is feasible since the pure BNE we assumed exists is a solution to it. Denote the 
resulting expected payoffs by $v^{\mathrm{row}}_{ij}$ and $v^\mathrm{col}_{ij}$. 

The solution thus obtained is a $(3\eps/4)$-BNE\@. To see this, we actually observe that  the equilibrium is well supported: for each type $i$, 
each strategy $j$ with $X^{\mathrm{row}}_{ij}\neq 0$ the payoff 
\[
v^{\mathrm{row}}_{ij}\ge M_i^{\mathrm{row}}-\frac{3\eps}{8} \ge \max_{j'}
v^{\mathrm{row}}_{ij'}- \frac{5\eps}{8}\,.
\] 

To complete the proof we now obtain a set of pure strategies $s_{\mathrm{row}}':\Ta_{\mathrm{row}}\ra S_{\mathrm{row}}$, $s_\mathrm{col}':\Ta_\mathrm{col}\ra S_\mathrm{col}$ by sampling 
$j=s_{\mathrm{row}}'(i)$ according to the probability distribution $X^{\mathrm{row}}_{ij}$. 

If $k=\Omega((\ln (nk))/\eps^2)$, then once again by the Chernoff-Hoeffding
bound the payoffs will be $\eps/8$-close to the payoffs
$v^{\mathrm{row}}_{ij}$, and thus the resulting game will be in an
$\eps$-BNE\@.  On the other hand if $k=O((\ln (nk))/\eps^2)$, \ie, the
number of types is small, we can find the exact BNE by brute force,
completing the proof of the theorem.
\end{proof}

%
%
%
%
%
\providecommand{\bibhead}[1]{}
%
\expandafter\ifx\csname pdfbookmark\endcsname\relax%
  \providecommand{\tocrefpdfbookmark}{}
\else\providecommand{\tocrefpdfbookmark}{%
   \hypertarget{tocreferences}{}%
   \pdfbookmark[1]{References}{tocreferences}}%
\fi

\tocrefpdfbookmark
\begin{thebibliography}{10}

\bibitem{AKS98}\bibhead{AKS98}
{\sc Noga Alon, Michael Krivelevich, and Benny Sudakov}: Finding a large hidden
  clique in a random graph.
\newblock {\em Random Structures {\&} Algorithms}, 13(3-4):457--466, 1998.
\newblock Preliminary version in
  \href{http://dl.acm.org/citation.cfm?id=314613.315014}{SODA'98}.
\newblock
  [\epfmtdoi{10.1002/(SICI)1098-2418(199810/12)13:3/4<457::AID-RSA14>3.0.CO;2-W}]

\bibitem{ABC11}\bibhead{ABC11}
{\sc Per Austrin, Mark Braverman, and Eden Chlamt\'{a}\v{c}}: Inapproximability
  of {NP}-complete variants of {N}ash equilibrium.
\newblock In {\em Proc. 14th Internat. Workshop on Approximation Algorithms for
  Combinatorial Optimization Problems (APPROX'11)}, pp. 13--25. Springer, 2011.
\newblock [\epfmtdoi{10.1007/978-3-642-22935-0\_2}]

\bibitem{BBM10}\bibhead{BBM10}
{\sc Hartwig Bosse, Jaroslaw Byrka, and Evangelos Markakis}: New algorithms for
  approximate {N}ash equilibria in bimatrix games.
\newblock {\em Theoret. Comput. Sci.}, 411(1):164--173, 2010.
\newblock Preliminary version in
  \href{http://dx.doi.org/10.1007/978-3-540-77105-0_6}{WINE'07}.
\newblock [\epfmtdoi{10.1016/j.tcs.2009.09.023}]

\bibitem{BV09}\bibhead{BV09}
{\sc S.~Charles Brubaker and Santosh Vempala}: Random tensors and planted
  cliques.
\newblock In {\em Proc. 13th Internat. Workshop on Randomization and
  Computation (RANDOM'09)}, pp. 406--419, 2009.
\newblock [\epfmtdoi{10.1007/978-3-642-03685-9\_31}]

\bibitem{CDT09}\bibhead{CDT09}
{\sc Xi~Chen, Xiaotie Deng, and Shang-Hua Teng}: Settling the complexity of
  computing two-player {N}ash equilibria.
\newblock {\em J. ACM}, 56(3):14:1--14:57, 2009.
\newblock Preliminary versions in FOCS'06 by
  \href{http://dx.doi.org/10.1109/FOCS.2006.69}{Chen and Deng} and by
  \href{http://dx.doi.org/10.1109/FOCS.2006.20}{Chen, Deng, and Teng}.
\newblock [\epfmtdoi{10.1145/1516512.1516516}]

\bibitem{CS03}\bibhead{CS03}
{\sc Vincent Conitzer and Tuomas Sandholm}: Complexity results about {N}ash
  equilibria.
\newblock In {\em Proc. 18th Internat. Joint Conf. on Artificial Intelligence
  (IJCAI'03)}, pp. 765--771, 2003.
\newblock Article accessible at
  \href{http://ijcai.org/Past\%20Proceedings/IJCAI-2003/PDF/111.pdf}{IJCAI}.
\newblock [\epfmt{acm}{1630659.1630770}]

\bibitem{CS08}\bibhead{CS08}
{\sc Vincent Conitzer and Tuomas Sandholm}: New complexity results about {N}ash
  equilibria.
\newblock {\em Games and Economic Behavior}, 63(2):621--641, 2008.
\newblock [\epfmtdoi{10.1016/j.geb.2008.02.015}]

\bibitem{dailey1980uniqueness}\bibhead{dailey1980uniqueness}
{\sc David~P. Dailey}: {Uniqueness of colorability and colorability of planar
  4-regular graphs are {NP}-complete}.
\newblock {\em Discrete Mathematics}, 30(3):289--293, 1980.
\newblock [\epfmtdoi{10.1016/0012-365X(80)90236-8}]

\bibitem{DMP07}\bibhead{DMP07}
{\sc Constantinos Daskalakis, Aranyak Mehta, and Christos~H. Papadimitriou}:
  Progress in approximate {N}ash equilibria.
\newblock In {\em Proc. 8th ACM Conf. on Electronic Commerce (EC'07)}, pp.
  355--358. ACM Press, 2007.
\newblock [\epfmtdoi{10.1145/1250910.1250962}]

\bibitem{DMP09}\bibhead{DMP09}
{\sc Constantinos Daskalakis, Aranyak Mehta, and Christos~H. Papadimitriou}: A
  note on approximate {N}ash equilibria.
\newblock {\em Theoret. Comput. Sci.}, 410(17):1581--1588, 2009.
\newblock Preliminary version in
  \href{http://dx.doi.org/10.1007/11944874_27}{WINE'06}.
\newblock [\epfmtdoi{10.1016/j.tcs.2008.12.031}]

\bibitem{FNS07}\bibhead{FNS07}
{\sc Tom\'{a}s Feder, Hamid Nazerzadeh, and Amin Saberi}: Approximating {N}ash
  equilibria using small-support strategies.
\newblock In {\em Proc. 8th ACM Conf. on Electronic Commerce (EC'07)}, pp.
  352--354, 2007.
\newblock [\epfmtdoi{10.1145/1250910.1250961}]

\bibitem{FK03}\bibhead{FK03}
{\sc Uriel Feige and Robert Krauthgamer}: The probable value of the
  {L}ov{\'a}sz--{S}chrijver relaxations for {M}aximum {I}ndependent {S}et.
\newblock {\em SIAM J. Comput.}, 32(2):345--370, 2003.
\newblock [\epfmtdoi{10.1137/S009753970240118X}]

\bibitem{FK08}\bibhead{FK08}
{\sc Alan~M. Frieze and Ravi Kannan}: A new approach to the planted clique
  problem.
\newblock In {\em IARCS Ann. Conf. on Foundations of Software Technology and
  Theoretical Computer Science (FSTTCS'08)}, pp. 187--198. Schloss Dagstuhl,
  2008.
\newblock [\epfmtdoi{10.4230/LIPIcs.FSTTCS.2008.1752}]

\bibitem{FudenbergTirole}\bibhead{FudenbergTirole}
{\sc Drew Fudenberg and Jean Tirole}: {\em Game Theory}.
\newblock MIT Press, 1991.

\bibitem{GZ89}\bibhead{GZ89}
{\sc Itzhak Gilboa and Eitan Zemel}: {N}ash and correlated equilibria: Some
  complexity considerations.
\newblock {\em Games and Economic Behavior}, 1(1):80--93, 1989.
\newblock [\epfmtdoi{10.1016/0899-8256(89)90006-7}]

\bibitem{HK11}\bibhead{HK11}
{\sc Elad Hazan and Robert Krauthgamer}: How hard is it to approximate the best
  {N}ash equilibrium?
\newblock {\em SIAM J. Comput.}, 40(1):79--91, 2011.
\newblock Preliminary version in
  \href{http://dl.acm.org/citation.cfm?id=1496849}{SODA'09}.
\newblock [\epfmtdoi{10.1137/090766991}]

\bibitem{Hoeff63}\bibhead{Hoeff63}
{\sc Wassily Hoeffding}: Probability inequalities for sums of bounded random
  variables.
\newblock {\em J. Amer. Statist. Assoc.}, 58(301):13--30, 1963.
\newblock [\epfmtdoi{10.1080/01621459.1963.10500830}]

\bibitem{LMM03}\bibhead{LMM03}
{\sc Richard~J. Lipton, Evangelos Markakis, and Aranyak Mehta}: Playing large
  games using simple strategies.
\newblock In {\em Proc. 4th ACM Conf. on Electronic Commerce (EC'03)}, pp.
  36--41. ACM Press, 2003.
\newblock [\epfmtdoi{10.1145/779928.779933}]

\bibitem{MV09}\bibhead{MV09}
{\sc Lorenz Minder and Dan Vilenchik}: Small clique detection and approximate
  {N}ash equilibria.
\newblock In {\em Proc. 13th Internat. Workshop on Randomization and
  Computation (RANDOM'09)}, pp. 673--685. Springer, 2009.
\newblock [\epfmtdoi{10.1007/978-3-642-03685-9\_50}]

\bibitem{Papadimitriou94}\bibhead{Papadimitriou94}
{\sc Christos~H. Papadimitriou}: On the complexity of the parity argument and
  other inefficient proofs of existence.
\newblock {\em J. Comput. System Sci.}, 48(3):498--532, 1994.
\newblock [\epfmtdoi{10.1016/S0022-0000(05)80063-7}]

\bibitem{TS08}\bibhead{TS08}
{\sc Haralampos Tsaknakis and Paul~G. Spirakis}: An optimization approach for
  approximate {N}ash equilibria.
\newblock {\em Internet Mathematics}, 5(4):365--382, 2008.
\newblock Preliminary version in
  \href{http://dx.doi.org/10.1007/978-3-540-77105-0_8}{WINE'07}.
\newblock [\epfmtdoi{10.1080/15427951.2008.10129172}]

\end{thebibliography}

